\chapter{Logical expressions}

Logical expressions can only have the value \texttt{.TRUE.} or \texttt{.FALSE.}. A logical expression can be formed by comparing arithmetic expressions using the following \emph{relational operators}:

\begin{lstlisting}[style=fortran77]
      .LT.  meaning <
      .LE.  meaning <=
      .GT.  meaning >
      .GE.  meaning >=
      .EQ.  meaning =
      .NE.  meaning /=
\end{lstlisting}

So you \emph{cannot} use symbols like $<$ or $>$ for comparison in Fortran 77, but you have to use the correct two-letter abbreviation enclosed by dots!

Logical expressions can be combined by the \emph{logical operators} \texttt{.AND. .OR. .NOT.} which have the obvious meaning.

\section{Logical variables and assignment}

Truth values can be stored in \emph{logical variables}. The assignment is analogous to the arithmetic assignment. Example:

\begin{lstlisting}[style=fortran77]
      logical a, b
      a = .TRUE.
      b = a .AND. 3 .LT. 5/2
\end{lstlisting}

The order of precedence is important, as the last example shows. The rule is that arithmetic expressions are evaluated first, then relational operators, and finally logical operators. Hence b will be assigned \texttt{.FALSE.} in the example above. Among the logical operators the precedence (in the absence of parenthesis) is that \texttt{.NOT.} is done first, then \texttt{.AND.}, then \texttt{.OR.} is done last.

Logical variables are seldom used in Fortran. But logical expressions are frequently used in conditional statements like the \texttt{if} statement.
